\documentclass[11pt,a4paper]{article}

\usepackage{amsmath,amssymb}
\usepackage[utf8]{inputenc}
\usepackage{hyperref}
\usepackage{geometry}
\usepackage{booktabs}
\usepackage{xcolor}
\usepackage{enumitem}
\usepackage{setspace}

\geometry{margin=1in}
\onehalfspacing

\newcommand{\RH}{\ensuremath{\mathrm{RH}}}

\title{%
  Two Axioms Remain:\\
  An Invitation to the Next Generation%
}
\author{
  Jason Robert Gochanour
}
\date{April 19, 2026}

\begin{document}

\maketitle

\begin{center}
\emph{For the student who hasn't started yet.}
\end{center}

\vspace{1cm}

% ================================================================
\section*{What This Is}
% ================================================================

We built something called the Cathedral: a 78-file computer program
that proves, with mathematical certainty, that the Riemann Hypothesis
is \emph{exactly equivalent} to the convergence of a specific
approximation problem. Every step is verified by a computer. Every
assumption is named. Every failed attempt is preserved.

The proof is complete---except for seven assumptions, called
\emph{axioms}. We couldn't prove them. We named them, classified
them, and documented exactly what they say and why they're hard.

This paper is an invitation. Those seven axioms are now your problem.

% ================================================================
\section*{Who Built This}
% ================================================================

A computer scientist in New Mexico. Not a professional mathematician.
No PhD in mathematics. No university affiliation. No research grant.

The tools: a laptop, an internet connection, Lean~4 (free, open
source), Mathlib (free, open source), and AI assistance (commercially
available).

The timeline: 23 days.

If you have a laptop and curiosity, you have everything that was
needed to build this. The barrier to entry for machine-verified
mathematics is now approximately zero.

% ================================================================
\section*{What You Need to Know}
% ================================================================

\subsection*{The Problem}

The Riemann Hypothesis (1859) says: all non-trivial zeros of the
Riemann zeta function $\zeta(s) = \sum_{n=1}^\infty n^{-s}$ lie
on the line $\operatorname{Re}(s) = \frac{1}{2}$.

Nobody has proved this. Nobody has disproved it. It has been open
for 167 years. There is a \$1,000,000 prize.

\subsection*{The Translation}

The Cathedral proves that this is equivalent to:

\[
  \forall\,\varepsilon > 0,\;\exists\,N_0,\;\forall\,N \geq N_0,\;
  \exists\,v \in \mathbb{R}^{N-1}:\quad
  \int_0^1 \left(1 - \sum_{k=2}^N v_k\left\{\frac{1}{kx}\right\}
  \right)^2 dx < \varepsilon.
\]

In words: you can approximate the constant function $1$ arbitrarily
well by mixing together sawtooth waves $\{1/(kx)\}$.

\subsection*{The Two Axioms}

\begin{enumerate}
  \item \textbf{\texttt{bd\_mellin\_at\_zero}}: The Mellin transform
    of the fractional-part function $\{1/(kx)\}$, which is proved
    for $\operatorname{Re}(s) > 1$ by direct computation, extends
    by analytic continuation to $\operatorname{Re}(s) > 0$.

    \emph{What this means}: a specific integral formula, proved for
    large values of a parameter, also holds for smaller values.
    This is a standard result in complex analysis, but formalizing
    it in Lean requires Mathlib's analytic continuation infrastructure.

  \item \textbf{\texttt{rh\_implies\_l2\_convergence}}: If the
    Riemann Hypothesis is true, then the B\'aez-Duarte distance
    $d_N^2 \to 0$.

    \emph{What this means}: assuming RH, you can construct explicit
    weights $v_k$ that make the integral above as small as you want.
    The construction uses the Mertens function, Abel summation, and
    the Montgomery--Vaughan mean value theorem.
\end{enumerate}

Axiom~1 is a formalization challenge: the mathematics is understood,
but translating it into Lean~4 requires careful work with analytic
continuation.

Axiom~2 is a \emph{mathematics} challenge: it packages several
deep results in analytic number theory. Proving it formally would
require formalizing the Mertens--Abel--Montgomery--Vaughan chain.

% ================================================================
\section*{What You Can Do Right Now}
% ================================================================

\subsection*{If You Know Some Lean~4}

\begin{enumerate}
  \item \textbf{Prove Axiom~1.} This is the more accessible target.
    You need Mathlib's \texttt{DifferentiableOn.analyticContinuation}
    or equivalent, applied to the Mellin integral
    $\int_0^1 \{1/(kx)\}\,x^{s-1}\,dx$. The integral is absolutely
    convergent for $\operatorname{Re}(s) > 0$; the challenge is
    showing that the closed-form expression (involving $\zeta(s)$
    and $1/(s-1)$) extends analytically.

  \item \textbf{Prove one of the 38 non-critical axioms.} The
    Cathedral contains 40 axioms total. Only 2 are on the crown
    path. The other 38 support alternative proof engines (spectral,
    sieve, Vasyunin). Proving any one of them strengthens the
    overall formalization.

  \item \textbf{Improve an existing proof.} Some of the Cathedral's
    644 theorems use \texttt{omega} or \texttt{norm\_num} where
    a more elegant proof exists. Professional Lean hackers will
    see opportunities to improve proof clarity and performance.
\end{enumerate}

\subsection*{If You Know Some Analysis}

\begin{enumerate}
  \item \textbf{Decompose Axiom~2.} The current axiom is a
    ``black box'' that packages the entire forward direction.
    A useful contribution would be to decompose it into 3--5
    intermediate lemmas, each of which is independently
    meaningful and potentially formalizable.

  \item \textbf{Find a better forward proof.} The B\'aez-Duarte
    theorem uses the Mertens function bound $|M(x)| = O(x^{1/2}
    \log^2 x)$ under RH. Is there a more direct path from RH
    to $d_N^2 \to 0$ that avoids the Mertens detour?

  \item \textbf{Study the decay rate.} The Cathedral's numerical
    experiments show $d_N^2 \sim c/\ln N$ with $c \approx
    1/21.65$. Can you prove this asymptotic? Can you prove a
    weaker bound, like $d_N^2 = o(1)$, from RH with fewer tools?
\end{enumerate}

\subsection*{If You Know Some Physics}

\begin{enumerate}
  \item \textbf{Identify the operator.} The Gram matrix $G_N$ is
    the finite-dimensional shadow of some infinite-dimensional
    operator whose spectral properties encode RH. What is this
    operator? The eigenvalue data ($\lambda_{\min} \sim C/\ln N$,
    $\lambda_{\max} \sim N$, GUE-like statistics) constrain it.

  \item \textbf{Interpret the constant.} The quantum stiffness
    $C \approx 21.65 = 1/(2 + \gamma - \ln 4\pi)$ appears to
    be a physical constant of the ``prime number vacuum.'' Does
    it arise naturally from any known quantum system?
\end{enumerate}

% ================================================================
\section*{What Modern Math Research Actually Looks Like}
% ================================================================

It doesn't look like a genius working alone in a garret.

It looks like this:

\begin{enumerate}
  \item \textbf{Day 1--7}: Explore. Try things. Fail. Archive
    the failures. We produced 96 files of failed attempts. Every
    failure taught us something.

  \item \textbf{Day 8--14}: Find the right framework. For us,
    this was the Parseval Bridge---the realization that frequency-
    domain methods bypass the hardest discrete computations. This
    came from a failure: the False Dedekind Reciprocity at
    $(a,b) = (3,2)$.

  \item \textbf{Day 15--21}: Build. Once the framework is right,
    construction is fast. AI handles the mechanical code generation.
    The human provides architecture. The compiler verifies
    everything.

  \item \textbf{Day 22--23}: Audit. Check every file. Fix every
    docstring. Remove every ghost axiom. Reduce, simplify, clean.
    The Great Audit reduced 178 files to 78 active + 96 archived.
\end{enumerate}

The ratio of exploration to construction is roughly 2:1. Two weeks
of wandering, one week of building. This is normal. If you're not
failing most of the time, you're not exploring hard enough.

% ================================================================
\section*{Why You Should Care}
% ================================================================

The Riemann Hypothesis is not just a mathematical curiosity. It
controls the distribution of prime numbers, which are the atoms
of arithmetic. Every encryption algorithm, every hash function,
every digital signature depends on properties of primes.

But more than that: the RH is a test case for whether humans can
understand the deepest structures in mathematics. If we can prove
it, we learn something profound about the nature of number. If we
can't, we learn something profound about the limits of mathematical
reasoning.

Either way, the attempt is worth making.

% ================================================================
\section*{A Note on Credentials}
% ================================================================

You do not need a PhD to contribute to this project. You do not
need a university appointment. You do not need a research grant.

You need:
\begin{itemize}
  \item A laptop.
  \item Lean~4 (free).
  \item Mathlib (free).
  \item Curiosity.
  \item The willingness to fail 96 times before you succeed once.
\end{itemize}

The compiler does not check your credentials. It checks your proof.

% ================================================================
\section*{The Torch}
% ================================================================

We built the Cathedral. We mapped the territory. We named the
axioms. We preserved the failures. We verified every step.

We could not prove the last two things.

Now it's your turn.

\begin{quote}
\emph{The primes have been singing for 167 years.\\
The Cathedral stands, waiting for its final stones.\\
We built it as far as we could.\\
The rest belongs to you.}
\end{quote}

\begin{thebibliography}{99}

\bibitem{baezduarte2003}
L.~B\'aez-Duarte, \emph{A strengthening of the Nyman--Beurling
criterion for the Riemann hypothesis}, Atti Accad. Naz. Lincei
Rend. Lincei Mat. Appl. \textbf{14} (2003), 5--11.

\bibitem{lean4}
L.~de~Moura \emph{et al.}, \emph{The Lean~4 Theorem Prover and
Programming Language}, CADE-28, 2021.

\bibitem{avigad}
J.~Avigad \emph{et al.}, \emph{Mathematics in Lean}, online
textbook, 2023.

\bibitem{buzzard}
K.~Buzzard, \emph{The Natural Number Game},
\url{https://adam.math.hhu.de/\#/g/leanprover-community/nng4}.

\end{thebibliography}

\end{document}
